\documentclass[thesis]{cluu}

\usepackage[style=cluu]{biblatex}
\addbibresource{main.bib}

\begin{document}
\title{My thesis}
\author{Xiaotian Luo}
% \date{...}
\supervisors{My Supervisor, Uppsala University\\
  Some Other Supervisor, NLP Enterprises, Inc.}
% Use \supervisor (without s at end) if you have only one

\maketitle

\begin{abstract}
    text here
\end{abstract}

\tableofcontents

\addchap{Preface}
% with \addchap instead of \chapter it isn't numbered.

\chapter{Introduction}
text here

\chapter{Background}
\section{Domain Adaptation of Dense Retrieval}
\subsection{From Lexical to Dense Retrieval}

Information retrieval (IR) is the task of identifying and ranking 
documents from a large corpus that are relevant to a given query. 
Early IR systems approached this through lexical methods, representing 
both queries and documents as sparse bag-of-words vectors. Models such 
as TF-IDF and BM25 compute relevance scores based on term frequency 
statistics and inverse document frequency weighting 
\parencite{manning2008introduction, robertson2009probabilistic}. By 
leveraging inverted indices, these methods enable efficient retrieval 
at scale and remain strong baselines to this day.

However, lexical retrieval suffers from a fundamental limitation known 
as \textit{vocabulary mismatch}: when a query and a relevant document 
share the same meaning but express it through different words, 
term-based methods fail to establish a match. This brittleness is 
particularly problematic in specialised domains where terminology is 
inconsistent or evolves over time.

Dense retrieval addresses this limitation by mapping queries and 
documents into a shared continuous vector space using neural networks, 
enabling matching based on semantic similarity rather than exact term 
overlap. Building on the contextual representations introduced by BERT 
\parencite{devlin-etal-2019-bert}, Sentence-BERT demonstrated that 
transformer-based models could be fine-tuned to produce meaningful 
sentence-level embeddings \parencite{reimers2019sentence}. This gave 
rise to the \textit{bi-encoder} architecture, in which a query and a 
document are encoded independently into dense vectors, and relevance 
is computed as their dot product or cosine similarity 
\parencite{karpukhin2020dense}. Since document embeddings can be 
pre-computed offline, approximate nearest neighbour search libraries 
such as FAISS \parencite{johnson2019billion} make dense retrieval 
computationally feasible even over millions of documents.

The strong performance of dense retrieval models depends critically 
on large-scale supervised training data consisting of 
(query, document) pairs. Datasets such as MS~MARCO, 
containing hundreds of thousands of human-annotated relevance 
judgements, have been central to training state-of-the-art dense 
retrievers \parencite{bajaj2016ms}. As we discuss in the following 
section, this dependency on labelled data becomes a significant 
obstacle when moving beyond general-domain web search.

\subsection{Domain Adaptation and the Pain of Label Scarcity}

The effectiveness of dense retrieval models comes at a cost: they 
require large quantities of labelled (query, document) pairs 
to learn meaningful vector representations. Outside of general-domain 
web search, such annotations are rarely available. In specialised 
fields such as medicine, law, or historical archives, producing 
relevance judgements requires significant domain expertise and is 
too expensive to do at scale. Models must therefore often be applied in 
a \textit{zero-shot} setting, where no labelled data from the target 
domain exists.

This creates a serious problem. Research using the BEIR benchmark, a 
diverse collection of retrieval datasets spanning biomedical, 
scientific, and other specialised domains, has shown that dense 
retrieval models trained on MS~MARCO perform poorly when applied to 
unseen domains \parencite{thakur2021beir}. In many cases, these 
neural models are outperformed by the simple lexical baseline BM25. 
The reason is that a bi-encoder learns a mapping function that places 
text into specific positions in vector space. When the target domain 
differs significantly from the training domain, this mapping no longer 
places relevant documents and queries close together. This gap between 
the training distribution and the target distribution is known as 
\textit{domain shift}.

Addressing domain shift without labelled data requires 
\textit{unsupervised domain adaptation} (UDA). The core idea is to 
generate synthetic training data directly from the target corpus, 
which can then be used to fine-tune the model on the new domain. 
Recent approaches such as GPL \parencite{wang-etal-2022-gpl} follow 
this direction and have shown promising results. We discuss these 
methods in detail in the following section.

\section{Related Work}

The field of unsupervised domain adaptation for dense retrieval has 
evolved through several distinct paradigms, each addressing the 
limitations of its predecessor.

\paragraph{\textbf{Domain-Adaptive Pre-training}}served as the initial 
approach, focusing on helping models internalize the vocabulary and 
linguistic patterns of the target corpus. Foundational methods like 
Masked Language Modeling (MLM) \parencite{devlin-etal-2019-bert} 
established general language understanding, while architectures such 
as Condenser \parencite{gao-callan-2021-condenser} and TSDAE 
\parencite{wang-etal-2021-tsdae} designed pre-training objectives 
specifically to improve dense retrieval performance. A significant limitation of these methods is that learning the 
vocabulary and patterns of a domain is not the same as learning 
to retrieve. Pre-training on target domain text teaches the model 
to better represent that text, but it provides no signal about 
which documents are relevant to a given query. The core skill of 
retrieval, matching a query to a relevant passage remains untrained.

\paragraph{\textbf{Self-Supervised Contrastive Learning}} attempted to solve 
this by creating artificial retrieval tasks without manual labels. 
For instance, the Inverse Cloze Task (ICT) \parencite{lee-etal-2019-latent} 
treats a random sentence as a query for its context, while SimCSE 
\parencite{gao-etal-2021-simcse} and Contrastive Tension (CT) 
\parencite{carlsson2021semantic} use dropout noise or dual encoders 
to generate contrastive pairs for fine-tuning. While these methods 
are efficient for learning sentence embeddings, they do not reflect 
real user search behaviour. Randomly extracted sentences from a 
document often lack the specific structure and information-seeking 
nature of a real-world user query.

\paragraph{\textbf{Pure Query Generation}} addressed this disconnect by using 
Large Language Models (LLMs) to synthesize human-like queries from 
target documents. Seminal works like Doc2Query 
\parencite{nogueira2019doc2query} and domain-targeted QGen 
\parencite{ma-etal-2021-zero} showed that fine-tuning on these 
synthetic pairs can bridge the document-to-query gap. Nevertheless, 
LLM-generated queries are inherently noisy and may contain 
hallucinations. This forces the model to treat all synthetic pairs 
as equally reliable, regardless of their actual quality.

\paragraph{\textbf{Generative Pseudo-Labeling (GPL)}} 
\parencite{wang-etal-2022-gpl} represents the current state-of-the-art 
solution and serves as the methodological framework for this thesis. 
GPL follows a three-step pipeline: first, an LLM generates synthetic 
queries for each passage in the target corpus; second, hard negatives 
are mined using a dense retriever; third, a Cross-Encoder scores each 
(query, passage) pair to produce soft relevance labels. By supervising 
the retriever with these continuous scores rather than binary labels, 
GPL effectively handles the noise inherent in synthetic query 
generation and produces more robust domain-adapted models. While more 
recent approaches such as DRAGON \parencite{lin-etal-2023-train} 
explore ensemble-based labeling to further improve training signal 
diversity, GPL remains the most widely adopted framework for 
unsupervised domain adaptation due to its practical balance between 
performance and computational cost. This thesis therefore adopts GPL 
as its methodological foundation.

\section{LLM-based Query Generation}
\begin{itemize}
  \item Data processing
  \item Prompt engineering
\end{itemize}

\section{Challenges in Historical Swedish Retrieval}
\begin{itemize}
  \item OCR errors
  \item Spelling variations
  \item Semantic shifts
  \item Boundary issues
\end{itemize}


\printbibliography
% at start yields errors because there are no citations and the bibliography file is empty

\end{document}
